% Resposta teórica em frequência
\begin{frame}{Resposta teórica em frequência}
  \begin{columns}[c,onlytextwidth]
    \begin{column}{0.58\textwidth}
      \centering
      \begin{tikzpicture}
        \begin{axis}[
          width=\linewidth,
          height=0.54\textheight,
          xmode=log,
          xmin=35,xmax=13000,
          ymin=-13,ymax=37,
          grid=both,
          major grid style={LavandaLinha},
          minor grid style={LavandaClaro},
          xlabel={Frequência (Hz)},
          ylabel={Magnitude (dB)},
          tick label style={font=\scriptsize},
          label style={font=\scriptsize},
        ]
          \addplot[very thick,RoxoPrincipal,mark=*,mark size=1.7pt]
          coordinates {
            (40,-4.51) (100,3.77) (200,11.05) (400,25.20)
            (480,33.97) (550,27.11) (1100,11.57) (2200,4.16)
            (5500,-4.15) (11000,-10.22)
          };
          \addplot[dashed,MalvaDestaque] coordinates {(478.9,-13) (478.9,37)};
        \end{axis}
      \end{tikzpicture}

      \scriptsize Curva teórica calculada a partir de $H(j\omega)$.
    \end{column}
    \begin{column}{0.38\textwidth}
      \begin{block}{Parâmetros centrais}
        \small
        \[
          f_0=\frac{1}{2\pi}
          \sqrt{\frac{R_1+R_2}{R_1R_2R_3C^2}}
          \approx 478{,}9\,\mathrm{Hz}
        \]
        \[
          Q\approx 7{,}07,
          \qquad |H(j\omega_0)|\approx 50.
        \]
      \end{block}
      \begin{alertblock}{Interpretação}
        O circuito atenua baixas e altas frequências e apresenta ganho máximo na
        vizinhança de $480\,\mathrm{Hz}$.
      \end{alertblock}
    \end{column}
  \end{columns}
\end{frame}